% ============================================================
% Capitolo 4 — _CONTEXT.md: la fonte di verità
% ============================================================

\chapter{\texttt{\_CONTEXT.md}: la fonte di verità}
\label{ch:context-md}

In un progetto software tradizionale, il contesto vive nelle teste delle persone. Con gli agenti AI il problema è amplificato: ogni sessione è un ``ingresso di qualcuno di nuovo''. L'agente non ha memoria. Senza un documento strutturato, sei tu che ricostruisci il contesto a voce ogni volta.

\texttt{\_CONTEXT.md} è la risposta a questo problema. È la fonte di verità del progetto: un singolo file Markdown che l'agente legge all'inizio di ogni sessione e che contiene tutto ciò che serve per lavorare correttamente.

\section{Dove vive e come viene trovato}
\label{sec:04-dove-vive}

\texttt{\_CONTEXT.md} vive alla radice del tuo progetto. L'agente segue questo protocollo di ricerca:

\begin{enumerate}
  \item Cerca \texttt{\_CONTEXT.md} nella directory corrente.
  \item Se non lo trova, risale nelle directory padre fino alla radice del repository.
  \item Se ne trova più di uno (monorepo), usa quello più vicino al file su cui sta lavorando.
  \item Se non trova nulla, chiede di crearne uno.
\end{enumerate}

\section{La struttura del file}
\label{sec:04-struttura}

Un \texttt{\_CONTEXT.md} completo ha quattro sezioni principali.

\subsection*{Sezione 1 --- Context card (obbligatoria)}

\begin{lstlisting}[language=,caption={Context card di TaskFlow API}]
| Param                    | Value                                       |
|--------------------------|---------------------------------------------|
| Project                  | TaskFlow API                                |
| Phase                    | 3-Implementation                            |
| Mode                     | STANDARD                                    |
| Active Task              | T-007.2 GET /tasks con filtri               |
| Active Task Token Est.   | 5000/1500/6500                              |
| Active Task Model Level  | 3 Sonnet Low                                |
| Active Branch            | feat/T-007-tasks-endpoints                  |
| Last Checkpoint          | 2026-06-27 10:47 — POST /tasks completato   |
\end{lstlisting}

\textbf{Phase} (0--6) determina quali moduli il framework carica. \textbf{Mode} controlla il livello di cerimonia (Capitolo~\ref{ch:fasi-e-modes}). \textbf{Active Task Token Est.} è la stima \texttt{input/output/totale} in token. \textbf{Active Task Model Level} (1--7) è il livello AI necessario.

\subsection*{Sezione 2 --- Stack (obbligatoria)}

\begin{lstlisting}[language=,caption={Stack tecnologico nel context}]
## Stack
| Layer      | Technology                  |
|------------|-----------------------------|
| Runtime    | Node.js 22                  |
| Framework  | Fastify 5                   |
| Database   | PostgreSQL 16 + Prisma ORM  |
| Auth       | OIDC + JWT (jose library)   |
| Testing    | Vitest + Supertest          |
| Deploy     | Docker + Railway            |
\end{lstlisting}

L'agente usa questa tabella per fare scelte coerenti senza dover chiedere.

\subsection*{Sezione 3 --- Vincoli di sicurezza (obbligatoria se applicabile)}

\begin{lstlisting}[language=,caption={Vincoli SEC nel context}]
## Security Constraints
| ID     | Name              | Spec                                |
|--------|-------------------|-------------------------------------|
| SEC-01 | Input Validation  | Zod schema obbligatorio             |
| SEC-02 | Authentication    | OIDC, JWT con rotazione             |
| SEC-03 | Logging           | Nessun token o PII nei log          |
\end{lstlisting}

\subsection*{Sezione 4 --- Vincoli di performance (obbligatoria se applicabile)}

\begin{lstlisting}[language=,caption={Vincoli PERF nel context}]
## Performance Constraints
| ID      | Name         | Target              |
|---------|--------------|---------------------|
| PERF-01 | Latency      | P95 < 200ms         |
| PERF-02 | Payload Size | Max 1MB per risposta|
\end{lstlisting}

\section{Sezioni opzionali utili}
\label{sec:04-sezioni-opzionali}

\subsection*{Decisioni architetturali aperte}

\begin{lstlisting}[language=,caption={Open Decisions nel context}]
## Open Decisions
| ID    | Decision                  | Options         | Deadline   |
|-------|---------------------------|-----------------|------------|
| AD-01 | Paginazione cursor/offset | cursor / offset | 2026-07-05 |
\end{lstlisting}

Se una decisione è qui, l'agente non la prende da solo --- ti chiede di decidere.

\subsection*{Note per l'agente}

\begin{lstlisting}[language=,caption={Note nel context}]
## Agent Notes
- Priorità: qualità > velocità
- Test: integration test su DB reale, no mock di Prisma
- Naming: snake_case DB, camelCase TS, kebab-case file
\end{lstlisting}

\section{Come aggiornare \texttt{\_CONTEXT.md}}
\label{sec:04-aggiornamento}

Aggiorna il file ogni volta che cambia:
\begin{itemize}
  \item La fase del progetto
  \item Il task attivo
  \item I vincoli SEC o PERF
  \item Una decisione aperta viene presa
\end{itemize}

Il modo più pratico è il comando \texttt{@context-update} (Capitolo~\ref{ch:comandi-conversazionali}).

\section{Template minimo vs template completo}
\label{sec:04-template}

Il framework offre due template in \texttt{.ai-dlc/modules/templates/}:

\begin{itemize}
  \item \textbf{CONTEXT\_TEMPLATE.md} --- completo, con tutte le sezioni incluse decisioni aperte e note.
  \item \textbf{CONTEXT\_MIN.md} --- minimo per spike e POC: solo context card, stack e vincoli essenziali.
\end{itemize}

\begin{tipbox}
Inizia sempre dal minimo. Un \texttt{\_CONTEXT.md} parzialmente compilato è meglio di uno vuoto. Aggiungi sezioni quando ne senti la necessità.
\end{tipbox}

\section{Errori comuni}
\label{sec:04-errori}

\textbf{Lasciare i placeholder.} Il validator segnala warning se \texttt{Active Task Token Est.} contiene \texttt{0/0/0}. Non blocca il lavoro, ma indica contesto incompleto.

\textbf{Non aggiornare il file.} Un \texttt{\_CONTEXT.md} non aggiornato è un contesto che mente. Dopo una settimana senza aggiornamenti, l'agente lavora su informazioni obsolete.

\textbf{Mettere troppo nel context.} Le regole dettagliate del progetto vanno in \texttt{.ai-dlc/project/instructions.md}. Il context contiene lo stato attuale, non il manuale del progetto.

\textbf{Non committare il file.} \texttt{\_CONTEXT.md} va in git. È parte del progetto, non un file temporaneo.

\section{Il contesto non è magia}
\label{sec:04-non-magia}

Vale la pena essere onesti su un punto che la ricerca recente ha messo in evidenza. Gli studi empirici sui file di contesto (\texttt{\_CONTEXT.md}, \texttt{AGENTS.md}, \texttt{CLAUDE.md}) danno risultati misti: un file mal progettato, troppo lungo o non aggiornato può aumentare il costo in token e perfino peggiorare l'aderenza dell'agente invece di migliorarla. Un \texttt{\_CONTEXT.md} non è un talismano --- è una scelta di design che vale solo se trattata con disciplina ingegneristica.

Tre principi che la letteratura conferma:

\begin{itemize}
  \item \textbf{Minimalità.} Più contesto non è meglio. Includi solo ciò che cambia il comportamento dell'agente; rimuovi il resto. Le panoramiche generiche del progetto sono poco utili --- i vincoli concreti, i comandi esatti e le modalità di fallimento note sono ciò che fa la differenza.
  \item \textbf{Aggiornamento.} Il fenomeno del \textit{context rot} --- l'invecchiamento del contesto rispetto al codice reale --- è documentato in letteratura: un \texttt{\_CONTEXT.md} che diverge dal repository diventa una fonte di errori, non di verità.
  \item \textbf{Specificità.} Un vincolo vago (\texttt{usa buone pratiche di sicurezza}) non orienta l'agente. Un vincolo specifico (\texttt{SEC-03: redaction di password, token, email nei log con Pino}) sì.
\end{itemize}

Il valore di \texttt{\_CONTEXT.md} non è automatico: dipende interamente da quanto bene lo mantieni. La classificazione del rischio e i confidence tag (Parte III) e il ciclo di checkpoint (Capitolo~\ref{ch:sessione-reale}) esistono proprio per imporre questa disciplina.

\section*{\LabelSummary}

\begin{itemize}
  \item \texttt{\_CONTEXT.md} è il singolo file che dà all'agente il contesto strutturato per lavorare correttamente senza istruzioni ripetute.
  \item Contiene: context card, stack, vincoli SEC e PERF, e sezioni opzionali.
  \item Va aggiornato a ogni cambio di fase, task o vincolo.
  \item Inizia dal template minimo; aggiungi sezioni man mano che il progetto cresce.
  \item Committalo in git e non lasciarlo mai stale.
\end{itemize}
